\section{Integral lattices and classical data}
\label{app:integral-charges}

\subsection{The resolved chamber and intersection form}

Adjoin $v_9=p=(-2,-1,0,0)$ and $v_{10}=(-1,0,0,0)$ to
\eqref{eq:x9-polytope}. The lower-facet triangulation with heights
\begin{equation}
 h_i=b_i+(i+1)/1009,\qquad
 (b_0,\ldots,b_{10})=(10,11,3,9,2,7,1,4,8,-5,-3)
\end{equation}
fixes a complete smooth fan with $29$ maximal cones. The divisor
$D_{10}$ misses the hypersurface. In the basis
\eqref{eq:x9-divisor-basis}, with $D_9=E$, the nonzero cubic
entries are as follows; all permutations agree and all omitted
entries vanish:
\begin{equation}
\begin{array}{cc|cc|cc|cc}
 IJK&\kappa_{IJK}&IJK&\kappa_{IJK}&IJK&\kappa_{IJK}&IJK&\kappa_{IJK}\\\hline
 222&8&266&1&667&-1&788&-6\\
 225&-2&267&1&677&-1&888&8\\
 226&-3&569&1&699&-2&889&-1\\
 227&-2&599&-2&777&-3&899&-1\\
 256&1&666&-1&778&4&999&7
\end{array}
\label{eq:x9-full-cubic-data}
\end{equation}
The second-Chern row is $(-4,24,26,18,-4,-2)$, and the
remaining divisor classes are
\begin{equation}
\begin{aligned}
 D_0&=D_2+4D_5+4D_6-3D_7-2D_8+2E,\\
 D_1&=D_2+2D_5+2D_6-D_7-D_8+E,\\
 D_3&=D_5-D_7-D_8,\qquad D_4=D_6-D_7-D_8.
\end{aligned}
\end{equation}
These data determine every norm, charge pairing and cubic used
in Sections~\ref{subsec:x9}--\ref{sec:decoupling}.

\subsection{Resolved and smoothed lattices}
Order the resolved divisors as $(D_2,D_5,D_6,D_7,D_8,E)$, with
$D_9=E$. The six integral curve classes
\begin{equation}
 (D_6D_7,D_2D_4,D_4D_5,D_8E,D_3D_7,D_0D_5)|_{\wideX_9}
\end{equation}
give the divisor--curve pairing, with curves as rows,
\begin{equation}
 P=\begin{pmatrix}
 1&0&-1&-1&0&0\\-1&1&0&1&0&0\\1&0&0&0&0&1\\
 0&0&0&0&-1&-1\\0&0&1&-1&2&0\\2&0&3&0&0&0
 \end{pmatrix},\qquad \det P=1.
 \label{eq:x9-integral-pairing}
\end{equation}
Since the six divisors are already a rational basis, any integral
cohomology class has rational coordinates whose pairing with these
curves is integral. Multiplication by $P^{-1}$ shows that those
coordinates are integral. Thus the displayed divisors generate
$H^2(\wideX_9;\ZZ)/\mathrm{tors}$, and the displayed curve classes
generate its dual. This argument does not infer integral saturation
from the rational Hodge-number count.

The Mayer--Vietoris description of the smooth fiber also holds on free
integral lattices. The links have $H^1(L;\ZZ)=0$. On the $E_2$ surface
use $(h,e_1,e_2)$ with intersection form $\operatorname{diag}(1,-1,-1)$
and $K=-3h+e_1+e_2$. For $V_7=\operatorname{Bl}_{pt}\mathbb P^3$,
\begin{equation}
 H|_E=2h-e_1-e_2,\qquad Q|_E=h-e_1-e_2,\qquad
 K=-2H|_E+Q|_E.
 \label{eq:x9-integral-local-gluing}
\end{equation}
The two restrictions form an integral basis of $R^\perp$, where
$R=e_1-e_2$: their matrix on $(h,e_1+e_2)$ has determinant $-1$.
Both $K$ and the deleted divisor class $2H-Q$ are primitive. Therefore
$H^2(V_7\setminus E;\ZZ)\cong\ZZ^2/\ZZ(2,-1)$ maps isomorphically
to $R^\perp/\ZZ K$ in the link cohomology, with no finite index.
For the resolution, excision and Poincar\'e--Lefschetz duality give
$H^3(\wideX_9,Y;\ZZ)\cong H_3(N;\ZZ)=0$, since the exceptional
neighborhoods retract onto $E$ and the two nodal curves. Thus restriction
to the common complement is surjective, with kernel $\ZZ E$; the nodal
Milnor fibers have $H^2=0$. The class $E$ is primitive by
\eqref{eq:x9-integral-pairing}. The neutral
lattice has integral basis
\begin{equation}
 (A,B,D_8,E)=(D_2+D_7,D_5+D_6-D_7,D_8,E).
\end{equation}
Indeed adjoining $-D_6,D_5$ gives a unimodular six-divisor basis.
Quotienting by $E$ gives \eqref{eq:x9-primitive-smooth-basis}.

\subsection{The local correction to the cubic}

Set $V_7=\operatorname{Bl}_{pt}\mathbb P^3$, with pullback
hyperplane $H$ and exceptional plane $Q$. The compactifying
surface is $S_7\in|2H-Q|$ and its complement is the Milnor
fiber. The restrictions in \eqref{eq:x9-integral-local-gluing}
identify the surviving link classes with $R^\perp/\ZZ K_E$.
Every inherited divisor in $\operatorname{span}(D_0,D_1)$
restricts to zero on the exceptional surface and the links.
It therefore has a compactly supported representative on $Y$;
all products with an inherited factor agree on both sides, as
do its second-Chern pairings.

For the remaining correction, glue the Milnor fiber to the
orientation-reversed exceptional neighborhood. As a cut and paste,
the resulting closed six-manifold is
$V_7-\mathbb P(\mathcal O\oplus K_{S_7}^{-1})$.
If $\lambda=D_8|_E=e$ is the exceptional curve on
$S_7=\operatorname{Bl}_1(\mathbb P^1\times\mathbb P^1)$,
its lifts are $Q$ on $V_7$ and $\pi^*\lambda$ on the projective
bundle. They give
\begin{equation}
 Q^3=1,\qquad (\pi^*\lambda)^3=0,\qquad
 p_1(V_7)Q=4,\qquad
 p_1\bigl(\mathbb P(\mathcal O\oplus K^{-1})\bigr)\pi^*\lambda=0.
\end{equation}
Here $c_1(V_7)=4H-2Q$, $c_2(V_7)H=6$ and $c_2(V_7)Q=0$.
Since $p_1=-2c_2$ on a Calabi--Yau threefold, these equalities
give $\Delta(\nu^3)=1$ and $\Delta(c_2\nu)=-2$.
The independence of the chosen resolved representative is explicit:
\begin{equation}
 (D_8+tE)^3+\bigl(-2tH+(1+t)Q\bigr)^3=9.
\end{equation}
The second term is the local correction associated with the same
boundary class, not an extra bulk divisor.

Geometrically, the $(-1)$-curve $D_8\cap E$ contracts and the
surface $D_8\cong\mathbb F_1$ maps to a plane through the
singularity. Its local cone over $e$ is replaced in the smoothing
by $Q\setminus e\cong\CC^2$. Gluing along their common
three-sphere gives a smooth plane of class $\nu$.
Adjunction then yields $\nu^3=K_{\mathbb P^2}^2=9$ and
$c_2(X_{9,t})\nu=c_2(\mathbb P^2)-9=-6$ independently.

\subsection{The integral flavor action}
The effective action on the reduced Higgs cone need not be faithful
on all states. In particular the $E_2$ non-Abelian flavor factor is
$SU(2)$, not $SO(3)$ \cite[Section~3.9.1]{Apruzzi:2021vcu}.
We check the compact action on the geometric M2 charge lattice,
using the structure-group criterion of
\cite[Section~2.1]{Apruzzi:2021vcu}: divide the product of gauge and
flavor factors by the central subgroup acting trivially on those charges.

For $\gamma=d h+m_1e_1+m_2e_2$, choose integral charges
\begin{equation}
 q_g=-K\cdot\gamma=3d+m_1+m_2,\qquad
 w=-R\cdot\gamma=m_1-m_2,\qquad q_a=h\cdot\gamma=d.
 \label{eq:x9-local-integral-charges}
\end{equation}
Here $q_g$ is the local Coulomb charge, $w$ is the standard $SU(2)$
weight with root weight two, and $q_a$ is a complementary Abelian
charge convention, allowing shifts by the Coulomb generator. Their
image is the index-two lattice $q_g+w+q_a\in2\ZZ$. Hence the
faithful structure group on this charge lattice is
\begin{equation}
 \frac{U(1)_g\times SU(2)_F\times U(1)_a}
 {\langle(-1,-\mathbf1,-1)\rangle}.
 \label{eq:x9-local-structure-group}
\end{equation}
This specifies the geometric charge action, not additional ultraviolet
flavor data or an independent identification of the six-dimensional
$U(2)$ gauge factor.

The restrictions of the six compact divisors to $E$, as columns in
$(h,e_1,e_2)$, are
\begin{equation}
 T=\begin{pmatrix}
 0&1&1&0&1&-3\\0&-1&0&0&-1&1\\0&0&-1&0&-1&1
 \end{pmatrix}.
 \label{eq:x9-full-restriction}
\end{equation}
For $D|_E=xh+ye_1+ze_2$, its action on a curve is
\begin{equation}
 D\cdot\gamma
 =-\frac{y+z}{2}q_g+\frac{z-y}{2}w
  +\left(x+\frac32(y+z)\right)q_a.
 \label{eq:x9-full-local-action}
\end{equation}
In particular, the coefficient triples for the two effective vectors are
\begin{equation}
 V_1:\ (-\tfrac12,\tfrac12,\tfrac12),\qquad
 V_2:\ (\tfrac12,\tfrac12,-\tfrac12).
 \label{eq:x9-cocharacter-lifts}
\end{equation}
Each $2\pi$ loop closes by the central identification in
\eqref{eq:x9-local-structure-group}. On the reduced current ring the
Coulomb and complementary Abelian actions are trivial, leaving unit
charges for $J_+$ under both vectors, as required in
\eqref{eq:x9-operator-charges}. Thus the reduced action is consistent
without requiring a half-Cartan to be a cocharacter of $SU(2)$ alone.

Finally this marking makes explicit the nilpotent-coordinate freedom
in \eqref{eq:e2-dictionary}. Let $U$ be the current normalized for
$q_a$. On the six compact divisors the coefficient of that current is
$(0,-\tfrac12,-\tfrac12,0,-2,0)$, whereas the geometric
$\alpha$ column is $(0,-2,1,0,-2,0)$. Their difference is
$\tfrac32$ times the root column. The compatible invariant coordinates are
\begin{equation}
 \alpha_{\mathrm{geom}}=U,\qquad
 \beta_{\mathrm{geom}}=J_0+\tfrac32U.
 \label{eq:x9-nilpotent-marking}
\end{equation}
The relations $U^2=UJ_0=0$ show directly that this shear preserves the
full invariant scheme. On its reduction $U=0$ it changes neither the
cooperative kernel nor the dressed-instanton ring.

\subsection{The formal invariant ring}
\label{subsec:x9-dressed-instantons}

In the primitive effective basis $(-D_6,D_5)$, the charges are
\begin{equation}
\begin{array}{c|rrrrrrr}
 &h_1&\widetilde h_1&h_2&\widetilde h_2&J_+&J_-&\beta\\\hline
 V_1&1&-1&0&0&1&-1&0\\
 V_2&0&0&1&-1&1&-1&0 .
\end{array}
\label{eq:x9-operator-charges}
\end{equation}
Remove the neutral pairs $h_i\widetilde h_i$ and $J_+J_-$
from any invariant monomial. Charge neutrality leaves a power
of $h_1h_2J_-$ or $\widetilde h_1\widetilde h_2J_+$.
Together with $u_i=-\beta$ and $J_+J_-=\beta^2$, this proves
the presentation \eqref{eq:x9-dressed-generators}. Grading the
hypermultiplet scalars by degree one and the currents by degree
two gives the refined Hilbert series
\begin{equation}
 H(t,z)=\frac{1-t^8}{(1-t^2)(1-zt^4)(1-z^{-1}t^4)}
       =\frac1{1-t^2}\sum_{n\in\ZZ}z^nt^{4|n|},
 \label{eq:x9-dressed-hilbert}
\end{equation}
where $Y$ has instanton fugacity $z$. An instanton current of
charge $n$ must be dressed by $|n|$ fields from each node.
Removing the local instanton currents forces $\beta^2=0$,
so the reduced formal quotient becomes a point. This detects
the local strong-coupling sector in the operator model; it is
not a heterotic perturbative calculation or a compact BPS index.

\subsection{Topological D-brane generators}
For either threefold $X$, let $L_a$ denote its ordered primitive divisor
basis: $H_1,\ldots,H_6$ on the resolution and $K_1,K_2,K_3$ on the
smooth side. Put $c_a=c_2(X)\cdot L_a$ and
$\kappa_{abc}=L_aL_bL_c$. Work in untwisted topological $K^0(X)$,
with Dirac pairing minus the Euler pairing. Choose generators
\begin{equation}
 \mathsf m_0=[\mathcal O_X],\qquad
 \mathsf m_a=[\mathcal O_X(L_a)]-[\mathcal O_X],\qquad
 \mathsf e_0=-[\mathcal O_{\rm pt}],\qquad
 \operatorname{ch}(\mathsf e_a)=(0,0,\beta_a,0),
 \label{eq:x9-d6-k-generators}
\end{equation}
where $L_a\cdot\beta_b=\delta_{ab}$. Sans-serif symbols denote charge
generators, distinct from the exceptional curves $e_1,e_2$.
The sign of $\mathsf e_0$ gives
normalized period $+1$, consistently with the earlier central-charge
convention. In particular $\mathsf m_0$ has Mukai charge
$(1,0,c_2(X)/24,0)$, not the unshifted cohomological unit.

These generators form the full integral charge lattice. On the smooth
side, the vertex matrices of the seven-ray fan polytope and its polar
both have nonzero Smith entries $(1,1,1,1)$. The same property was
established for the resolution. The fundamental-group and Brauer-group
formulas therefore give torsion-free integral cohomology on both sides
\cite[Corollaries~1.9 and~3.9]{Batyrev:2005Integral}.
The untwisted Atiyah--Hirzebruch spectral sequence degenerates: its
differentials are rationally zero and its groups are torsion-free.
Its free associated graded groups are $H^0,H^2,H^4,H^6$
\cite[Section~3.1]{Maldacena:2001KCharges}.
Thus every primitive $H^4$ class has a rank-zero, first-Chern-zero
$K$-theory lift. Its index is $\int_X\operatorname{ch}_3\in\ZZ$,
so a point class removes its third Chern character. This constructs
the $\mathsf e_a$ in \eqref{eq:x9-d6-k-generators}; the remaining generators
lift the primitive $H^0,H^2,H^6$ bases. No BPS representative of every
generator is assumed. In particular the smooth-side products of two
divisors generate an index-ten sublattice of $H^4$, not all of $H^4$.


\subsection{The current-exchange matrices}

For $J=J_*+\Lambda w$ with
$J_*=(6,18,16,-11,-8,7)$ in $\mathcal D$, both nodal areas
are one. The cubic tensor above and \eqref{eq:x9-full-inverse}
give
\begin{equation}
\begin{aligned}
 \mathsf H_{D_0}
 &=\frac1{66\Lambda}\begin{pmatrix}31&-9\\-9&31\end{pmatrix}
   +O(\Lambda^{-2}),\\
 \mathsf H_{D_1}
 &=\frac1{14}\begin{pmatrix}1&-1\\-1&1\end{pmatrix}
   +\frac1{196\Lambda}\begin{pmatrix}96&-5\\-5&110\end{pmatrix}
   +O(\Lambda^{-2}),\\
 \mathsf H_{D_0+D_1}
 &=\frac1{590\Lambda}\begin{pmatrix}209&-91\\-91&209\end{pmatrix}
   +O(\Lambda^{-2}).
\end{aligned}
\label{eq:x9-current-asymptotics}
\end{equation}
Since $Gt=v/2$ and $t^{\mathsf T}Gt=3\mathcal V$, graviphoton
subtraction removes the matrix
\[
 (3\mathcal V)^{-1}\begin{pmatrix}1&1\\1&1\end{pmatrix}
 =O(\Lambda^{-3}).
\]
